% SEGMENT OLP-0544-S01
% Part: set-theory
% Chapter: z
% Section: milestone

\documentclass[../../../include/open-logic-section]{subfiles}

\begin{document}

\olfileid{sth}{z}{milestone}
\olsection{$\Z^-$: ஒரு முக்கியக் கட்டம்}

அடுத்த பிரிவில் \stagesinf{} ஐ மீண்டும் கருதுவோம். எனினும், முடிவிலி
அடிகோளுடன் நாம் ஒரு முக்கியக் கட்டத்தை அடைந்துள்ளோம். $\Zminus$ என்ற
கோட்பாட்டிற்குத் தேவையான எல்லா அடிகோள்களும் இப்போது நம்மிடம் உள்ளன.
விரிவாக:

\begin{defn}
$\Zminus$ கோட்பாடு பின்வரும் அடிகோள்களைக் கொண்டுள்ளது: உறுப்புசார் சமத்துவம்,
சேர்ப்பு, ஜோடிகள், அடுக்குக்கணங்கள், முடிவிலி, மேலும் பிரிப்பு அடிகோள்
திட்டத்தின் எல்லா நேர்வுகளும்.
\end{defn}

இப்பெயர் (பின்னர் நாம் காணவிருக்கும் ஏதோவொன்று \emph{நீக்கப்பட்ட})
\emph{செர்மேலோ} கணக் கோட்பாட்டைக் குறிக்கிறது. செர்மேலோ இப்பெருமைக்குத்
தகுதியானவர்; ஏனெனில் அவர் தமது \citeyear{Zermelo1908Untersuchungen} ஆம்
ஆண்டுப் படைப்பில் அடிப்படையில் இக்கோட்பாட்டையே வகுத்தார்.\footnote{இதன்
வரலாறும் தொழில்நுட்ப விவரங்களும் பற்றிய சுவையான குறிப்புகளுக்கு
\citet[Appendix A]{Potter2004} ஐப் பார்க்கவும்.}

மிகப் பெருமளவு கணிதத்தைச் செய்வதற்கு இக்கோட்பாடு போதுமான ஆற்றலுடையது.
குறிப்பாக, நீங்கள் \olref[sfr][][]{part} முழுவதையும் திரும்பிப் பார்த்து,
அங்கு முறைசாராமல் செய்த அனைத்தையும் $\Zminus$ க்குள் இன்னும் முறைசார்ந்த
விதத்தில் செய்ய முடியும் என்று \emph{உறுதிப்படுத்திக்கொள்ள வேண்டும்}.
(அதைச் சிறிது நேரம் செய்தபின், முன்சென்று
\olref[sth][z][arbintersections]{sec} ஐ வாசிக்க விரும்பலாம்.) ஆகவே
இனிமேல், வேறு எந்தக் குறிப்புமின்றி, நாம் $\Zminus$ இல் (குறைந்தபட்சம்)
செயல்படுகிறோம் என எடுத்துக்கொள்வோம்.

\end{document}
